\section{Power graphs of every rank}
\label{sec:all-rank-graph-v158}
For $1\leq r\leq n$, let $P^{(r)}\subset\Pj(\Sym^2V)$ be the
rank-exactly-$r$ locus. Define $\mathcal G_{r,h}$ as the
scheme-theoretic closure of
\begin{equation}\label{eq:rank-graph-definition-v158}
 P^{(r)}\longrightarrow
 \Pj(\Sym^2V)\times\prod_{p=1}^{hr}\Pj(B_h(V)_p),
 \qquad[A]\longmapsto([A],[A],\ldots,[A^{hr}]).
\end{equation}
Keeping the final power is essential when $r<n$. Write $\mathcal U$
for the tautological rank-$r$ bundle on $\Gr(r,V)$ and
$\mathrm{CQ}(\mathcal U)$ for its relative complete-quadric space.
For rank one this means $\Pj(\Sym^2\mathcal U)=\Gr(1,V)$.

\begin{theorem}[The rank-stratified power graph]
\label{thm:all-rank-graph-v158}
There is a canonical isomorphism
\begin{equation}\label{eq:rank-graph-isom-v158}
 \mathcal G_{r,h}\simeq\mathrm{CQ}(\mathcal U)
                  \longrightarrow\Gr(r,V).
\end{equation}
The final power in \eqref{eq:rank-graph-definition-v158} recovers the
image $r$-plane by the degree-$2h$ Pluecker embedding. In particular
the graph is smooth and projective, independently of $h$. For its
relative boundary divisors $E_i$, $1\leq i<r$, its power base ideals
have total transforms
\begin{equation}\label{eq:rank-boundary-v158}
 J_p\OO_{\mathcal G_{r,h}}=
 \OO_{\mathcal G_{r,h}}\left(-\sum_{i=1}^{r-1}(p-hi)_+E_i\right),
 \qquad 0\leq p\leq hr.
\end{equation}
The projective power maps in this formula have their common divisors
removed. All powers $p>hr$ vanish on the rank-$r$ locus.
\end{theorem}
\begin{proof}
For $U\subset V$, the map $B_h(U)\to B_h(V)$ from the presentation
\eqref{eq:envelope-presentation-v158} is injective: any linear projection
$V\to U$ splitting the inclusion induces a retraction on the quotient
algebras. The injection itself is independent of that projection.
Locally this gives subbundles $B_h(\mathcal U)_p\subset B_h(V)_p$.
If $A$ has image $U$ and rank $r$, it is a nonsingular element of
$\Sym^2U$, and $A^{hr}$ generates the socle
$B_h(U)_{hr}=(\det U)^{2h}$. As $U$ varies these lines give the
homogeneous embedding
\[
 \Gr(r,V)\hookrightarrow\Pj(\mathbb S_{(2h)^r}V)
                         \hookrightarrow\Pj(B_h(V)_{hr}).
\]
Indeed the line at the coordinate $r$-plane is a highest-weight line
of weight $(2h)^r$; its orbit map is the Pluecker embedding followed
by its degree-$2h$ linear system. In particular this is a closed
immersion, not just a recovery of the image plane on closed points.

Over the recovered Grassmannian, the remaining graph is precisely the
relative simultaneous power graph for $B_h(\mathcal U)$. By
Theorem~\ref{thm:power-graph-v157} on each frame chart, it is
$\mathrm{CQ}(\mathcal U)$. The relative graph embedding, the subbundle
injections above, and the final-power Grassmannian embedding together
give a closed immersion into the product in
\eqref{eq:rank-graph-definition-v158}. Its dense open is the graph
of $P^{(r)}$. Both its source and that graph closure are integral,
so equality of their scheme-theoretic images proves
\eqref{eq:rank-graph-isom-v158} without a normalization step.
Applying \eqref{eq:power-boundary-v157} in rank $r$ proves the
transform formula. Local split inclusion of the coefficient bundles
ensures that ambient and relative power ideals agree. The vanishing
assertion follows from the socle degree of $B_h(U)$.
\end{proof}

The varieties of complete singular quadrics are classical: see
\cite[Definition~2.5, Remark~2.6 and Theorem~2.14]{CasarottiCornianiMassarenti2023}.
The point of the theorem is the exact recovery of their image-plane
projection and entire scheme structure from the power maps in one
fixed algebra $B_h(V)$, simultaneously for every $r$. The classical
orbit stratification is background, not an additional novelty claim.

\section{Contacts and minimal indices of arbitrary pencils}
\label{sec:singular-pencil-v158}
Let $R=\langle A_0,A_1\rangle\subset\Sym^2V$ be a genuine pencil,
$\Lambda=\Pj(R)$, and $r$ its normal rank. Its tautological symmetric
map is $A:V^*\otimes\OO_\Lambda\to V\otimes\OO_\Lambda(1)$.
Write
\[
 K=\ker A,\qquad U=K^\perp\subset V\otimes\OO_\Lambda.
\]
The kernel is saturated and hence a subbundle on the smooth curve;
$U$ is the saturation of the generic image. Thus $A$ is a generically
nonsingular section of $\Sym^2U\otimes\OO_\Lambda(1)$.
The map on the rank-$r$ open lifts uniquely to
$\widetilde\Lambda\to\mathcal G_{r,h}$ by properness. Its
Grassmannian projection is the map defined by $U$, and its source is
still $\Lambda$. Let $D_i$ be its contact divisor with $E_i$.

\begin{lemma}[Symmetric diagonalization with a kernel]
\label{lem:symmetric-dvr-v158}
Over $\C[[t]]$, a symmetric matrix of rank $r$ over the fraction field
is congruent to $\diag(t^{e_1},\ldots,t^{e_r},0,\ldots,0)$, with
$e_1\leq\cdots\leq e_r$. For the local matrix of a genuine pencil,
$e_1=0$. The exponents are its nonzero Smith exponents.
\end{lemma}
\begin{proof}
Choose a matrix entry of smallest valuation. If a diagonal entry
has that valuation, it is a pivot. Otherwise an off-diagonal entry
does, and the value of the quadratic form on $e_i+c e_j$ has the
same valuation for some $c\in\C$: its leading coefficient is a
nonzero polynomial in $c$, since $2$ is invertible. Extend this
primitive vector to a basis. The pivot divides every entry in its
row, so symmetric row-and-column elimination clears the row and
column over the ring. All remaining entries have valuation at least
that of the pivot. Repeat until the remaining block is zero.
Units have square roots in $\C[[t]]$ and can be absorbed into the
basis. Invertible congruence preserves determinantal ideals, which
identify the resulting exponents with Smith exponents. A genuine
pencil has a nonzero matrix at every projective parameter, giving a
unit entry and $e_1=0$. This proof supplies the local congruence step
rather than relying on left-right Smith form alone; compare the
classical pencil congruence treatment in \cite{Thompson1991}.
\end{proof}

For $j\geq0$ define the finite multiplication-syzygy map
\begin{equation}\label{eq:toeplitz-map-v158}
 C_j(R):V^*\otimes\Sym^jR^*\longrightarrow
                         V\otimes\Sym^{j+1}R^*,\qquad q\longmapsto A q,
 \qquad\kappa_j=\dim\ker C_j(R).
\end{equation}
Set $\kappa_{-1}=\kappa_{-2}=0$. Under a basis of $R$ this is the block
Toeplitz map with rows $A_0q_0$, $A_0q_i+A_1q_{i-1}$, and $A_1q_j$;
its definition is independent of that basis.

\begin{theorem}[Complete power-and-syzygy data]
\label{thm:singular-complete-data-v158}
For every complex symmetric pencil, including every singular pencil,
the following formulas recover its complete congruence data.
If $\varepsilon_1,\ldots,\varepsilon_c$ are its minimal indices,
where $c=n-r$, then
\begin{align}
 K&\simeq\bigoplus_{a=1}^c\OO_\Lambda(-\varepsilon_a),
 &\kappa_j&=\sum_{a=1}^c(j-\varepsilon_a+1)_+,
                       \label{eq:kernel-splitting-v158}\\
 \#\{a:\varepsilon_a=j\}
   &=\kappa_j-2\kappa_{j-1}+\kappa_{j-2}.
                       \label{eq:minimal-differences-v158}
\end{align}
It is enough to use $0\leq j\leq\lfloor(n-1)/2\rfloor$.
At $x\in\Lambda$, let $0=e_1(x)\leq\cdots\leq e_r(x)$ be the
nonzero local Smith exponents. With $v_0=0$ and
$v_p(x)=\length(\OO_{\Lambda,x}/J_p)$ for $p\leq hr$, one has
\begin{align}
 v_{hq+s}(x)&=h\sum_{i=1}^{q}e_i(x)+s e_{q+1}(x),
                         \label{eq:singular-valuations-v158}\\
 e_i(x)&=\frac{v_{hi}(x)-v_{h(i-1)}(x)}h,
 &\operatorname{mult}_xD_i&=e_{i+1}(x)-e_i(x).
                         \label{eq:singular-contacts-v158}
\end{align}
In the first formula $s=0$ at $q=r$ and the last term is omitted.
The positive exponents at each supported point, together with the
minimal indices, determine the pencil up to congruence and projective
reparametrization of $\Lambda$. The spectral datum includes the
positions of the points, not only their local partitions.

There is also a global conservation law:
\begin{equation}\label{eq:conservation-v158}
 \sum_{i=1}^{r-1}(r-i)\deg D_i
       +2\deg\bigl(\Lambda\longrightarrow\Gr(r,V)\bigr)
       =r,\qquad
 \deg\bigl(\Lambda\longrightarrow\Gr(r,V)\bigr)
       =\sum_a\varepsilon_a.
\end{equation}
The Grassmannian degree uses its Pluecker polarization. The first
summand is the dimension of the regular summand in the Kronecker
congruence form. Every object and number in the theorem is invariant
under the ungraded failure-algebra isomorphisms of
Theorem~\ref{thm:intrinsic-envelope-v158}.
\end{theorem}
\begin{proof}
We use the classical complex symmetric Kronecker congruence theorem
\cite{Thompson1991}. A singular block of minimal index $\varepsilon$
has size $2\varepsilon+1$ and form
\[
 S_\varepsilon(s,t)=
 \begin{pmatrix}0&L_\varepsilon(s,t)^{\mathsf t}\\
                 L_\varepsilon(s,t)&0\end{pmatrix},\qquad
 L_\varepsilon=
 \begin{pmatrix}s&t& &0\\ &s&t& \\ &&\ddots&t\end{pmatrix},
\]
where $L_\varepsilon$ has $\varepsilon$ rows and $\varepsilon+1$
columns; $S_0$ is a zero $1\times1$ block. Its kernel is generated
by $((-t)^\varepsilon,s(-t)^{\varepsilon-1},\ldots,s^\varepsilon)$
in the first block. These components have no common projective zero,
so it gives $\OO(-\varepsilon)$ as a kernel subbundle. A regular
block has zero kernel as a sheaf. Direct sums prove the first formula
in \eqref{eq:kernel-splitting-v158}. Taking sections after twist by
$\OO(j)$ identifies them with \eqref{eq:toeplitz-map-v158}, proving
the second formula. The second finite difference of
$(j-\varepsilon+1)_+$ is one exactly at $j=\varepsilon$ and zero
otherwise. Each singular block uses $2\varepsilon+1\leq n$
dimensions, proving the stated finite bound. The canonical-form
classification is used here as a classical input, not reproved or
claimed as new.

Locally, symmetric diagonalization factors the matrix through the
saturated $r$-plane $U$ with an $r\times r$ symmetric block. Theorem
\ref{thm:universal-power-v157}, or its rank-$r$ version via the split
inclusion $B_h(U)\hookrightarrow B_h(V)$, yields
\eqref{eq:singular-valuations-v158}. Finite differences recover the
$e_i$. In the relative complete-quadric chart the successive ratios
of diagonal entries have valuations $e_{i+1}-e_i$, proving the contact
formula. The torsion of the pencil cokernel has precisely the positive
Smith exponents. The regular blocks of the classical congruence form
are specified by these elementary divisors, and its singular blocks
by the minimal indices; over $\C$ there are no additional signature
parameters. Retaining the supported points on $\Lambda$ makes the
statement compatible with arbitrary projective changes of the pencil
basis.

The exact sequence $0\to U\to V\otimes\OO\to K^*\to0$ gives
$\deg U=-\sum\varepsilon_a$, so the Grassmannian degree is
$\sum\varepsilon_a$. The determinant of the induced symmetric map
$U^*\to U(1)$ is a nonzero section of
$(\det U)^2\otimes\OO(r)$; its zero divisor has degree
$r-2\sum\varepsilon_a$. At each point its order is
$\sum e_i=\sum_{i<r}(r-i)\operatorname{mult}_xD_i$.
This proves \eqref{eq:conservation-v158}. The block-size identity
$n=n_{\mathrm{reg}}+\sum(2\varepsilon_a+1)$ identifies the first
summand with $n_{\mathrm{reg}}$. Finally the intrinsic first relation
extracts $\Lambda_A$ on its source and preserves all the maps involved;
changes of source lift or pencil basis give isomorphic kernel maps
and equal ideals, hence the asserted invariance.
\end{proof}

\begin{example}[Power contacts alone miss a singular invariant]
\label{ex:singular-separation-v158}
The size-six pencils $S_0\oplus S_2$ and $S_1\oplus S_1$ both have
constant rank four and no spectral torsion. All their power ideals
are the unit ideal for $p\leq4h$ and zero for $p>4h$, and every
contact divisor is empty. Their kernel splittings are respectively
$\OO\oplus\OO(-2)$ and $\OO(-1)^2$. Thus their first three
syzygy nullities are $(1,2,4)$ and $(0,2,4)$, distinguishing the two
congruence types. Both have Grassmannian degree two, as required by
\eqref{eq:conservation-v158}; the full splitting, rather than just
its degree, is necessary.
\end{example}

\begin{proposition}[Finite closed tests in families]
\label{prop:syzygy-family-v158}
For $j\geq0$ and $0\leq a\leq n(j+1)$, on the pencil
Grassmannian each condition $\kappa_j\geq a$ is the
closed determinantal condition on the bundle map $C_j$ given by minors
of size $n(j+1)-a+1$. These conditions commute with arbitrary base
change as schemes. Their fibrewise ranks need not be constant, so no
unqualified assertion that the kernel bundles commute with every
base change is made. Along a specialization over an integral curve,
$\kappa_j$ can only increase.
\end{proposition}
\begin{proof}
Use the tautological rank-two subbundle in
\eqref{eq:toeplitz-map-v158}. Its source has rank $n(j+1)$ and its
target rank $n(j+2)$. Nullity at least $a$ is equivalent to rank at
most $n(j+1)-a$, giving the stated minors. The ideals are polynomial
in the bundle map and therefore commute with arbitrary substitution.
Upper semicontinuity of nullity, or specialization of these minors,
gives the last assertion.
\end{proof}
